\begin{figure}[!htbp]
\centering
\begin{adjustbox}{max width=\linewidth}
\begin{tikzpicture}
  \fill[fslate!35] (-1,2.5) rectangle (11,4.1); \fill[fsage!45] (-1,0.9) rectangle (11,2.3); \fill[fsand!55] (-1,-0.7) rectangle (11,0.7);
  \node[capt, anchor=west] at (-0.9,3.85) {external level}; \node[capt, anchor=west] at (-0.9,2.05) {conceptual level}; \node[capt, anchor=west] at (-0.9,0.45) {internal level};
  \foreach \x/\n in {3/1, 5.4/2, 7.8/3} {
    \node[box, fill=white, minimum width=20mm, font=\footnotesize] (v\n) at (\x,3.2) {external view \n};
    \node[font=\scriptsize] (u\n) at (\x,4.55) {user \n};
    \draw[arr] (u\n) -- (v\n);
  }
  \node[box, fill=white, text width=58mm, font=\footnotesize] (C) at (5.4,1.6) {\textbf{conceptual schema} — entities, relationships, constraints};
  \node[box, fill=white, text width=58mm, font=\footnotesize] (I) at (5.4,0.0) {\textbf{internal schema} — storage structures, indexes, files};
  \node[db, minimum width=26mm, minimum height=13mm] (D) at (5.4,-1.8) {stored\\database};
  \foreach \n in {1,2,3} \draw[arrd] (v\n) -- (C);
  \draw[arrd] (C) -- (I); \draw[arrd] (I) -- (D);
  \node[lbl, anchor=west] at (8.8,2.4) {logical data independence};
  \node[lbl, anchor=west] at (8.4,0.8) {physical data independence};
\end{tikzpicture}
\end{adjustbox}
\caption{The ANSI/SPARC three-schema architecture. Mappings between adjacent levels insulate each level from
changes in the one below it.}
\end{figure}
